\section{Periods and the discriminant on a selected deformation space}
\label{sec:marked-periods}

Write $B=\Def^S(X)$. At each degree-six vertex take the simultaneous
incidence modification on the chosen local component. Its central fibre
has two ordinary double points on the line component and three on the
plane component. All the original ordinary double points are retained.
Write $Z\to X$ for this modification on the central fibre and
$C_1,\ldots,C_n$ for the exceptional curves of the compatible small
resolution $Y\to Z$. The index of each curve records its original node
or its position over a marked cone. Orient the curves over a cone by
their complex orientations. At an original node a change of small
resolution changes the sign of its marked class. Put
\[
 R=\ker\!\left(\CC^n\longrightarrow H_2(Y;\CC),\quad
                         w\longmapsto\sum_iw_i[C_i]\right).
\]
The support of $w$ is the set of indices with $w_i\ne0$.

\begin{theorem}\label{thm:marked-periods}
The marked nodal periods define a holomorphic submersion of germs
$P\colon(B,0)\to(R,0)$. There are local analytic coordinates on $B$
in which $P$ is projection onto $R$. For each $i$ the ordinary double
point of the simultaneous modification persists exactly on
$P_i^{-1}(0)$, with its scheme structure as the pullback of the local
nodal discriminant. In particular every support occurring in $R$ occurs
on a projective deformation of $X$ inducing the selected profile.
\end{theorem}

\begin{proof}
Restrict the analytic family of $X$ over $B$ to representatives of the
marked germs. Pull back the simultaneous families of
\cite[\S3.4 and Proposition~3.8]{W5} along their classifying maps.
The modifications are proper, are flat over the parameter line or
plane, and are isomorphisms away from the vertex over the local
parameter origin. Thus the pullbacks glue to the identity on the
complement of smaller vertex neighborhoods. Properness is local on
the target, and flatness over $B$ follows from the two types of
charts. This gives a proper flat analytic family $g\colon\mathcal Z\to B$
with central fibre $Z$. A block whose local parameter is zero has the
central modification over that vertex, even on a fibre where another
vertex has moved. At a nonzero local parameter the map to the
corresponding part of the fibre of $X$ is an isomorphism, including
on the plane discriminant.

We use one relative volume form for all the nodes. Indeed, $Z$ is
connected and Gorenstein with $\omega_Z\simeq\cO_Z$ and
$H^1(Z,\cO_Z)=0$ by \cite[Proposition~4.1]{W5}. Serre duality
gives $H^2(Z,\cO_Z)=0$. After shrinking $B$, the relative dualizing
sheaf is invertible and commutes with base change. Semicontinuity
gives $h^0(\cO_{Z_b})=1$ and $h^1(\cO_{Z_b})=h^2(\cO_{Z_b})=0$.
Constancy of the holomorphic Euler characteristic in the proper flat
family gives $h^3(\cO_{Z_b})=1$. By duality and Grauert's base change
theorem a generator of $H^0(Z,\omega_Z)$ extends to a relative section
$\Omega$. Its zero locus misses the central fibre, so properness lets
us shrink $B$ until $\Omega$ generates $\omega_{\mathcal Z/B}$.

Remove mutually disjoint Milnor neighborhoods of the $n$ nodes of
$Z$. The compact exterior $W$ and its boundary are differentiably
locally trivial over $B$, by the proper submersion theorem with
boundary. The boundary component $L_i$ is diffeomorphic to
$S^2\times S^3$. Let $\gamma_i$ generate $H_3(L_i;\ZZ)$, oriented
so that the boundary of the Poincar\'e--Lefschetz dual on $W$ of
$D|_W$, for $D\in H^2(Y;\ZZ)$, has $i$th component
$\langle D,C_i\rangle\gamma_i$. This convention is possible since
restriction from a small-resolution neighborhood to the link is an
isomorphism on $H^2$. Transport the boundary cycles and set
$P_i(b)=\int_{\gamma_i}\Omega_b$. On the exterior $\Omega_b$ is
closed. Stokes' theorem on the transported relative four-cycle gives
\[
              \sum_i\langle D,C_i\rangle P_i(b)=0 .
\]
Integral degree-two classes span $H^2(Y;\CC)$ over $\CC$. The
evaluation pairing therefore gives $P(b)\in R$ on every fibre.

In a nodal chart write $\sum_{j=1}^4 x_j^2=\tau_i(b)$. The relative
form is $u(x,b)$ times the residue form, for $u$ a holomorphic unit.
When $\tau_i(b)\ne0$, inclusion of the link into the local Milnor
fibre induces an isomorphism on $H_3$: the fibre is a disc cotangent
bundle of $S^3$, with trivial sphere bundle as boundary. Up to the
chosen sign, $\gamma_i$ thus represents the sphere
$\sqrt{\tau_i(b)}\,S^3$. Scaling the residue form shows that its
integral is a fixed nonzero constant times $\tau_i(b)$ times the
spherical average of $u(\sqrt{\tau_i(b)}\,x,b)$. The odd Taylor
terms have zero average and the even terms converge normally on
a smaller chart. The average is holomorphic in $b$ and has nonzero
constant term $u(0,0)$. When $\tau_i(b)=0$, the cycle bounds in
the local small resolution and the form extends there by crepancy.
Consequently, also for an identically persistent node, we have
\begin{equation}\label{eq:period-unit}
               P_i(b)=\tau_i(b)v_i(b),\qquad v_i(0)\ne0
\end{equation}
with $v_i$ holomorphic. The units in this calculation come from the
same global form. Rescaling that form multiplies the whole vector
$P$ by a single unit on $B$.

Intersecting a relation with a vertex divisor $E_p$ gives
$\sum_{i\text{ over }p}w_i=0$: each of its marked curves has
$E_p$-degree $-1$, and all the other marked curves have degree zero.
The line block is consequently $(\lambda,-\lambda)$, and the plane
block can be written $(-a,b,a-b)$. The curve identifications of
\cite[Proposition~3.7 and Lemma~4.4]{W5} identify $R$ in these
coordinates with $K_S\otimes_{\QQ}\CC$. The local-to-global
tangent sequence in that lemma gives the image of the differential
of the selected local classifying map, of dimension $\dim R$.
Its subsequent differential to the tuple of nodal parameters is
injective on the local target: it is the identity at an original
node, records a nonzero multiple of the line parameter at a line
vertex, and on the plane records the three differences of the three
diagonal entries modulo simultaneous translation. These differences
separate tangent vectors in the plane. By \eqref{eq:period-unit},
multiplication by the diagonal matrix with entries $v_i(0)$ leaves
this rank unchanged. Since $P$ takes values in $R$, its differential
at zero is surjective.

Theorem~\ref{thm:regular} and the holomorphic submersion theorem
now give the asserted coordinates. Given $w\in R$, take a
local section through the origin and lift $t\mapsto tw$. Equation
\eqref{eq:period-unit} puts this entire disc in every prescribed
zero divisor and off every other one for $t\ne0$. A zero block is the
branch origin, since the nodal parameters jointly separate the
points of each local line or plane. A nonzero plane block with one
zero coordinate leaves precisely one ordinary double point. Outside
the marked neighborhoods the family is smooth after shrinking.

Let $f\colon\mathcal X\to\Delta$ denote the resulting disc family
of $X$, and fix an ample bundle $L$ on $X$. Topological proper base
change identifies $H^2(X;\ZZ)$ with the direct limit of
$H^2(f^{-1}(U);\ZZ)$ over discs $U$ about zero. Thus $c_1(L)$
extends after shrinking. Serre duality gives $H^2(X,\cO_X)=0$.
By semicontinuity, Grauert's theorem and the Stein property of the
disc, $H^2(\mathcal X,\cO_{\mathcal X})=0$ there. The exponential
sequence lifts the extended class to a line bundle, which restricts
to $L$ since $H^1(X,\cO_X)=0$. Openness of ampleness gives a
projective family even when $Z$ is nonprojective.
\end{proof}

For a projective conifold transition this arrangement is the one
described by Lee, Lin and Wang \cite[Lemma~3.13 and
Proposition~3.16]{LLW}. In the present construction $Z$ may itself
have a node which persists on the whole base. Such an index contributes
the zero functional on $R$. Two nonzero proportional coordinate
functionals give the same hyperplane, with both node markings retained.
The submersion concerns periods. The local equation parameters in
\eqref{eq:period-unit} may differ from them by different units.

\section{Cohomology of the regular locus and the marked torus}
\label{sec:marked-generators}

Let $I$ be a support in $R$. Denote by $F(I)$ the set of degree-six
vertices for which every index over the vertex lies outside $I$.
Put $H=H^2(Y;\ZZ)$, write $[E_p]\in H$ for the Poincar\'e dual
of the vertex divisor, and let
\[
 A_I=\{D\in H:\langle D,C_i\rangle=0\ \text{for all }i\in I\}.
\]
For a nearby fibre $X_b$ with $\operatorname{supp}P(b)=I$, the
following statement uses the markings transported from the central
fibre. Write $X_b^{\rm reg}$ for its smooth locus.

\begin{proposition}\label{prop:marked-generators}
There is an identification
\[
 H^2(X_b^{\rm reg};\ZZ)
       \simeq A_I\Big/\sum_{p\in F(I)}\ZZ[E_p].
\]
It identifies this group with a subgroup of
$H^2(X^{\rm reg};\ZZ)=H/\sum_p\ZZ[E_p]$. A class with lift $D\in H$
belongs to the image precisely when it has degree zero on every
original node in $I$ and its degrees on the curves in $I$ over each
one vertex are equal. For a representative $D\in A_I$, the restriction
at a residual ordinary double point is its degree on the marked
exceptional line. At a persistent cone $p$ it is
\[
 [D|_{E_p}]\in H^2(E_p;\ZZ)/\ZZ K_{E_p}\simeq H^2(L_p;\ZZ),
\]
where $L_p$ denotes the link of that cone.
\end{proposition}

\begin{proof}
Use the coordinates of Theorem~\ref{thm:marked-periods} to join the
origin to $b$ by an arc along which the support is $I$ away from the
origin. On a block in $F(I)$ the incidence family is the product of
the central model with that arc. Fill the exterior $W$ at the
indices in $I$ by the Milnor fibres, and at the remaining indices
by the compatible local small resolutions. Write $M_I$ for the
resulting smooth compact manifold. The filling which uses small
resolutions at all indices gives $Y$.

For either filling, each $H^1(L_i;\ZZ)$ is zero. On $H^2$ the
restriction of a small-resolution neighborhood to its boundary is an
isomorphism, and the Milnor fibre has $H^2=0$. Mayer--Vietoris first
identifies $H^2(Y;\ZZ)$ with $H^2(W;\ZZ)$ and then identifies the
image of $H^2(M_I;\ZZ)$ there with $A_I$. A vertex in $F(I)$
reconstructs the original smooth surface $E_p\subset M_I$, with
the marking of its divisor class preserved in the common exterior.
Remove these pairwise disjoint surfaces. By the Thom sequence,
the degree-two cohomology of the complement is obtained by quotienting
by their divisor classes, since $H^1(E_p;\ZZ)=0$. The exceptional
curves at all other retained nodes are disjoint from the removed
surfaces. Removing them has no effect in degree two, by the Thom
isomorphism in real codimension four. This complement is the smooth
locus of $X_b$. The case $I=\varnothing$ gives the central
identification in the statement.

The resulting map to $H^2(X^{\rm reg};\ZZ)$ is induced by
$A_I\hookrightarrow H$. A combination $\sum_pc_p[E_p]$ which
annihilates the curves in $I$ has $c_p=0$ for each vertex with an
index in $I$, as one sees by evaluating on any such curve. This
gives injectivity on the stated quotient. Changing a lift by
$c_p[E_p]$ subtracts the integer $c_p$ from each degree over $p$
and has no effect at other vertices or original nodes. The degrees
in a nonempty block can all be made zero precisely when they are
equal. At an ordinary double point restriction to the link agrees
with the degree on the exceptional line by the same filling calculation.
At a persistent cone the link is the circle bundle of $K_{E_p}$.
The Gysin sequence gives the stated quotient since $H^1(E_p;\ZZ)=0$;
adding $[E_p]$ to the lift adds $K_{E_p}$ to the restriction by adjunction.
This also describes the link classes outside the selected root summand.
\end{proof}

For the rigid abelian gauging, one can see both the supports and
the stabilizers directly. Work at zero triplet moment-map level with
just the marked fields. Use one charged hypermultiplet at each
original node and the presentation by two or three homogeneous
hypermultiplets at a chosen cone component, divided there by the
diagonal circle \cite[\S\S2--3]{Higgs}. The three moment maps for
one homogeneous hypermultiplet $(q_i,\widetilde q_i)$ comprise its
complex product $q_i\widetilde q_i$ and its real moment
$(|q_i|^2-|\widetilde q_i|^2)/2$. Their joint vanishing is equivalent
to $(q_i,\widetilde q_i)=0$. A common rotation of the three moment
maps makes the complex projection nonzero on all the nonzero
triplets of a given solution: the discarded real axis need only
avoid finitely many lines. The resulting complex products satisfy
the linear charge relations and give a vector of $R$ with exactly
that support. Conversely, for $w\in R$ take each product equal to
$w_i$ and both absolute values equal to $\sqrt{|w_i|}$. This solves
all three equations, including each diagonal-circle equation.
Together with Theorem~\ref{thm:marked-periods}, this gives equality
of the sets of supports realized on the two sides.

For the connected bulk torus take the cocharacter lattice
$H^2(X^{\rm reg};\ZZ)$ modulo torsion, with the original marked
degrees. Its action fixes the class of a configuration precisely
when a diagonal-circle action at each cone compensates the bulk
action on every nonzero homogeneous hypermultiplet there. On the
Lie algebra this requires equal degrees on the participating curves
over that cone and zero degree at each participating original node.
For a cocharacter with integral degrees the compensating common
degree is integral as well. Proposition~\ref{prop:marked-generators}
thus gives the marked lattice of the identity component of the
stabilizer, including spectator directions. An empty block leaves
its diagonal circle redundant, as accounted for by the quotient
by $[E_p]$. Finite components of a stabilizer depend in addition
on saturation of the subgroup of the character lattice generated
by the active weights.

With a conical Ricci-flat metric on the regular locus, its
degree-two $L^2$ cohomology identifies with middle-perversity
intersection cohomology \cite{Cheeger}. At an isolated point of
real codimension six, the middle-perversity truncation retains
the link cohomology through degree two. Thus
$IH^2(X_b;\mathbb R)=H^2(X_b^{\rm reg};\mathbb R)$. In a fluxless
three-form reduction these classes give the linearized normalizable
bulk vector fields. The local conicality criterion for bounded-potential
Ricci-flat metrics at isolated Calabi--Yau cones with smooth link is
given in \cite[Theorem~1.4]{Zhang}. The cohomological calculation
identifies the marked vector space on the geometric fibre. A
microscopic identification of the complete Higgs vacuum at finite
Planck mass would require further information about the localized
interacting theory.
